\chapter{Ausblick}

\begin{lernziele}
Nach diesem Kapitel kennen Sie sinnvolle nächste Schritte zur Vertiefung.
\end{lernziele}

\section{SysML vertiefen}
Nach dem Einstieg lohnt es sich, SysML genauer zu lernen. Besonders relevant sind saubere Modellorganisation, Wiederverwendung von Modellelementen, Variantenmodellierung und Schnittstellenkonzepte.

\section{SysML v2}
SysML v2 modernisiert die Sprache und bringt eine präzisere textuelle und grafische Modellierung. Für Einsteiger ist SysML v1.x weiterhin verbreitet, aber SysML v2 wird langfristig wichtiger.

\section{ROS2 vertiefen}
Für Robotik sind Navigation2, TF2, Lifecycle Nodes, Launch-Dateien, Parameter und Simulation wichtige nächste Themen.

\section{Simulation}
Gazebo oder andere Simulatoren ermöglichen Tests ohne echte Hardware. Das ist besonders wertvoll, wenn Algorithmen schnell ausprobiert werden sollen.

\section{Safety Engineering}
Sobald Roboter mit Menschen in realen Umgebungen interagieren, werden Sicherheitsfragen wichtig. Dann reichen einfache Funktionsmodelle nicht mehr aus. Risikoanalysen und Normen müssen berücksichtigt werden.

\section{Berufliche Relevanz}
Die Kombination aus Systems Engineering, SysML, ROS2 und KI ist in Robotik, Automotive, Industrieautomation und Forschung relevant. Wer diese Verbindung versteht, kann zwischen Architektur, Software und Tests vermitteln.

\begin{uebungbox}
Definieren Sie Ihren persönlichen nächsten Lernschritt: SysML vertiefen, ROS2 implementieren, Simulation aufbauen oder ML-Integration ausprobieren.
\end{uebungbox}
